% GENERATED FILE. Edit canonical lessons, then run npm run generate:books.
% canonical-source-hash: fnv1a64:18922240eac5cbaa
% canonical-lessons: GE-C12-elf, GE-C12-zwoelf, GE-C12-lif, GE-C12-elf-zwoelf

\chapter{One Left Over}
\label{ch:elf-zwoelf}
\hlchaptermodality{18}

\begin{chapteropening}
\textbf{By the end of this chapter:} \emph{I can say eleven and twelve in German and explain why neither of them follows the counting pattern.}

elf and zwoelf are not irregular numbers, they are a different rule, and it is a rule about hands. ainlif is one-left and twalif is two-left: left over from ten, because ten is where your fingers stop. English inherited both words unchanged, which is why eleven and twelve are odd there in exactly the same way -- when two languages share an irregularity this precisely they did not agree on it, they were one language when it was made. The rule covers two numbers and then stops.
\end{chapteropening}

\section[elf]{elf — \textquotedblleft{}eleven\textquotedblright{}}

\label{lesson:GE-C12-elf}



\begin{quote}
You can count to ten. Both languages then do something odd for two numbers, and they do the same odd thing.
\end{quote}

\subsection*{You'll want to know: elf}
\begin{quote}
\textbf{elf} — \textquotedblleft{}eleven.\textquotedblright{}
\end{quote}

Not \emph{einszehn}. Not \textquotedblleft{}oneteen\textquotedblright{} in English either. Eleven is the first of two numbers that refuse the pattern, in German and in English alike.

\begin{sounds}
\begin{itemize}
\raggedright
  \item \texttt{herz-short-e} — a short open \emph{e}, the vowel of English \emph{bed}.
  \item \texttt{f-final} — and a clean \emph{f} to close.
\end{itemize}

\emph{elf}, one syllable, said as it looks. This is the rare German word an English reader can pronounce on sight.
\end{sounds}

\subsection*{Your turn}
Say these aloud:

\begin{itemize}
\raggedright
  \item \textquotedblleft{}elf\textquotedblright{}
  \item the two in order — \textquotedblleft{}zehn, elf\textquotedblright{}
  \item what it is not — \textquotedblleft{}not einszehn\textquotedblright{}
\end{itemize}

\emph{Twice through:} \textquotedblleft{}zehn, elf.\textquotedblright{}

\begin{tcolorbox}[breakable,colback=teal!4,colframe=teal!35!black,title={Before you move on}]
Say \textquotedblleft{}eleven.\textquotedblright{} (\emph{\textbf{Elf}}.) Does it end in \emph{-zehn}? (\textbf{No}.) Does English say \textquotedblleft{}oneteen\textquotedblright{}? (\textbf{No} — \emph{eleven}.) Next: the second of the two.
\end{tcolorbox}
\section[zwölf]{zwölf — \textquotedblleft{}twelve\textquotedblright{}}

\label{lesson:GE-C12-zwoelf}



\begin{quote}
Short on the page, and three rules deep in the mouth. You own all three already.
\end{quote}

\subsection*{You'll want to know: zwölf}
\begin{quote}
\textbf{zwölf} — \textquotedblleft{}twelve.\textquotedblright{}
\end{quote}

The second number to refuse the pattern, and the last. After this one German turns regular and stays regular.

You can hear \emph{zwei} inside it if you listen for the opening: \emph{zw-} is the same start.

\begin{sounds}
\begin{itemize}
\raggedright
  \item \texttt{z-ts} — German \emph{z} is \textbf{ts}.
  \item \texttt{w-as-v} — German \emph{w} is a \textbf{v}, as on \emph{Winter} and \emph{Wasser}.
  \item \texttt{umlaut-oe} — and \emph{ö} is the rounded vowel from \emph{schön}: shape your lips for \emph{oh} and say \emph{eh}.
\end{itemize}

So the word opens \emph{tsv-} and runs \emph{tsvölf}. Three rules, one syllable, nothing new.
\end{sounds}

\subsection*{Your turn}
Say these aloud:

\begin{itemize}
\raggedright
  \item the opening alone — \textquotedblleft{}tsv\textquotedblright{}
  \item the vowel alone — lips for \textquotedblleft{}oh\textquotedblright{}, say \textquotedblleft{}eh\textquotedblright{}
  \item \textquotedblleft{}zwölf\textquotedblright{}
  \item the pair — \textquotedblleft{}elf, zwölf\textquotedblright{}
\end{itemize}

\emph{Twice through:} \textquotedblleft{}elf, zwölf.\textquotedblright{}

\begin{tcolorbox}[breakable,colback=teal!4,colframe=teal!35!black,title={Before you move on}]
Say \textquotedblleft{}twelve.\textquotedblright{} (\emph{\textbf{Zwölf}}.) How does it open? (\emph{\textbf{Tsv-}}.) Why? (\emph{\textbf{z}} \textbf{is} \emph{ts}, \emph{\textbf{w}} \textbf{is} \emph{v}.) Which number is hiding at the front? (\emph{\textbf{Zwei}}.) Next: what both of them actually mean.
\end{tcolorbox}
\section[-lif]{-lif — \textquotedblleft{}left over\textquotedblright{}}

\label{lesson:GE-C12-lif}



\begin{quote}
Two words that look like exceptions. They are not exceptions. They are an older rule, and it is a rule about hands.
\end{quote}

\begin{cousinweb}
Both words are compounds, and both end in the same piece:

\noindent
\begin{tabularx}{\linewidth}{@{}>{\raggedright\arraybackslash}X>{\raggedright\arraybackslash}X>{\raggedright\arraybackslash}X@{}}
\toprule
\textbf{word} & \textbf{older form} & \textbf{the pieces} \\
\midrule
\emph{elf} & \emph{ainlif} & \emph{ain} \textquotedblleft{}one\textquotedblright{} + \emph{lif} \\
\emph{zwölf} & \emph{twalif} & \emph{twa} \textquotedblleft{}two\textquotedblright{} + \emph{lif} \\
\bottomrule
\end{tabularx}

And \emph{\textbf{-lif}} means \textquotedblleft{}\textbf{to leave, to remain}.\textquotedblright{} So \emph{elf} is \textquotedblleft{}\textbf{one left over}\textquotedblright{} and \emph{zwölf} is \textquotedblleft{}\textbf{two left over}.\textquotedblright{}
\end{cousinweb}

\begin{culture}
Left over from \textbf{what}? From \textbf{ten} — from your ten fingers.

You count along them, you reach \emph{zehn}, you run out of hand, and there is one left. That is eleven. Not a number in a sequence but a \textbf{remainder}, named by people counting on the only calculator they had.

And English says the identical thing. \emph{Eleven} is \emph{ain-lif}; \emph{twelve} is \emph{twa-lif}. Not similar words — \textbf{the same two words}, inherited side by side, the way \emph{Vater} and \emph{father} are. The Germanic twinning you have been tracking through the household and the meal reaches all the way into the counting.

It also explains why the trick stops at two. Ten fingers leave room for one and two before the hands are hopeless, and after that a language needs a real rule.
\end{culture}

\subsection*{Your turn}
Say these aloud:

\begin{itemize}
\raggedright
  \item the two pieces of eleven — \textquotedblleft{}ain, lif\textquotedblright{}
  \item what the ending means — \textquotedblleft{}left over\textquotedblright{}
  \item the whole idea — \textquotedblleft{}one left over from ten\textquotedblright{}
  \item the English pair that says the same — \textquotedblleft{}eleven, twelve\textquotedblright{}
\end{itemize}

\emph{Twice through:} \textquotedblleft{}elf, zwölf — one left, two left.\textquotedblright{}

\begin{tcolorbox}[breakable,colback=teal!4,colframe=teal!35!black,title={Before you move on}]
What does \emph{-lif} mean? (\textbf{Left over}.) Left over from what? (\textbf{Ten} — \textbf{your fingers}.) So what is \emph{elf}? (\textbf{One left over}.) Are \emph{eleven} and \emph{twelve} the same words? (\textbf{Yes} — inherited, not borrowed.) Next: both of them together.
\end{tcolorbox}
\section[Practice]{Practice — one left, two left}

\label{lesson:GE-C12-elf-zwoelf}



\begin{quote}
No new words. Two numbers, and the reason they look strange in both languages at once.
\end{quote}

\begin{grammarlens}[title={Grammar Lens: what the two are made of}]
\noindent
\begin{tabularx}{\linewidth}{@{}>{\raggedright\arraybackslash}X>{\raggedright\arraybackslash}X>{\raggedright\arraybackslash}X@{}}
\toprule
\textbf{German} & \textbf{is} & \textbf{and means} \\
\midrule
\emph{elf} & \emph{ain} + \emph{lif} & one left over \\
\emph{zwölf} & \emph{twa} + \emph{lif} & two left over \\
\bottomrule
\end{tabularx}

Two words, one ending, and the ending is a verb: to leave, to remain. Neither number was ever irregular — they were built by a rule that only had room to run twice.
\end{grammarlens}

\begin{culture}
Left over from \emph{zehn}, from ten fingers. A counting system built on hands has to say something when the hands run out, and what Germanic said was \emph{and one is left}.

English inherited both words unchanged, which is why \emph{eleven} and \emph{twelve} are equally odd there, and odd in exactly the same way. When two languages share an irregularity this precisely, they did not agree on it — they were one language when it was made.

From the next number on, both of them turn regular and never break again.
\end{culture}

\subsection*{Your turn}
Say these aloud:

\begin{itemize}
\raggedright
  \item count the last three — \textquotedblleft{}zehn, elf, zwölf\textquotedblright{}
  \item the opening of the second — \textquotedblleft{}tsv\textquotedblright{}
  \item what each one means — \textquotedblleft{}one left over, two left over\textquotedblright{}
  \item the English pair — \textquotedblleft{}eleven, twelve\textquotedblright{}
\end{itemize}

\emph{Twice through:} \textquotedblleft{}zehn, elf, zwölf.\textquotedblright{}

\begin{tcolorbox}[breakable,colback=teal!4,colframe=teal!35!black,title={Before you move on}]
Say eleven and twelve. (\emph{\textbf{Elf, zwölf}}.) What does the ending mean? (\textbf{Left over}.) Left over from what? (\emph{\textbf{Zehn}} — \textbf{ten fingers}.) Why are the English words odd in the same way? (\textbf{They are the same two words}.) How many numbers does this rule cover? (\textbf{Two} — \textbf{and then it stops}.)
\end{tcolorbox}
